\documentclass[11pt]{article}

% ── Packages ────────────────────────────────────────────────────────────────
\usepackage[margin=1in]{geometry}
\usepackage{amsmath, amssymb, amsthm}
\usepackage{mathtools}
\usepackage{enumitem}
\usepackage{hyperref}
\usepackage{parskip}

% ── Theorem environments ─────────────────────────────────────────────────────
\newtheorem{theorem}{Theorem}[section]
\newtheorem{lemma}[theorem]{Lemma}
\newtheorem{corollary}[theorem]{Corollary}
\newtheorem{proposition}[theorem]{Proposition}
\theoremstyle{definition}
\newtheorem{definition}[theorem]{Definition}
\newtheorem{example}[theorem]{Example}
\theoremstyle{remark}
\newtheorem{remark}[theorem]{Remark}

% ── Custom commands ───────────────────────────────────────────────────────────
\newcommand{\R}{\mathbb{R}}
\newcommand{\E}{\mathbb{E}}
\newcommand{\norm}[1]{\left\lVert#1\right\rVert}
\newcommand{\inner}[2]{\langle #1,\, #2 \rangle}
\newcommand{\T}{^{\top}}

% ── Document ─────────────────────────────────────────────────────────────────
\begin{document}

\title{\textbf{Regression: Decision}\\[0.4em]
       \large Notes Template}
\author{}
\date{}
\maketitle
\tableofcontents
\newpage

% ─────────────────────────────────────────────────────────────────────────────
\section{Problem}
% ─────────────────────────────────────────────────────────────────────────────

\subsection{Setup}

% TODO: Describe the regression decision problem here.
% Example prompts:
%   - What is the input/output space?
%   - What loss function is used?
%   - What decision rule are we trying to derive?

\begin{definition}[Regression Decision Problem]
    % TODO: Fill in the formal definition.
    Let $\mathcal{X} \subseteq \R^{d}$ be the input space and
    $\mathcal{Y} \subseteq \R$ the output space.
    Given a learned model $\hat{f}: \mathcal{X} \to \mathcal{Y}$,
    the decision problem is to \ldots
\end{definition}

\subsection{Assumptions}

\begin{enumerate}[label=\textbf{A\arabic*.}]
    \item % TODO: State assumption 1.
    \item % TODO: State assumption 2.
    \item % TODO: Add further assumptions as needed.
\end{enumerate}

\subsection{Objective}

% TODO: State the objective (e.g., minimise expected loss, maximise posterior, etc.)
\[
    % TODO: Replace with the actual objective.
    \hat{y} \;=\; \arg\min_{y \in \mathcal{Y}} \; \mathcal{L}(y, \hat{f}(x))
\]

% ─────────────────────────────────────────────────────────────────────────────
\section{Derivation}
% ─────────────────────────────────────────────────────────────────────────────

The following derivation establishes the optimal decision rule step by step.

% ── Step 1 ────────────────────────────────────────────────────────────────────
\subsection{Step 1: [Title]}

% TODO: Describe what this step accomplishes.

\begin{proof}
    % TODO: Fill in the proof for Step 1.
    \[
        % placeholder
        \text{(derivation goes here)}
    \]
\end{proof}

% ── Step 2 ────────────────────────────────────────────────────────────────────
\subsection{Step 2: [Title]}

% TODO: Describe what this step accomplishes.

\begin{proof}
    % TODO: Fill in the proof for Step 2.
    \begin{align}
        % TODO: align environment for multi-line algebra.
        &\text{(left-hand side)} \notag \\
        &= \text{(right-hand side)} \label{eq:step2}
    \end{align}
\end{proof}

% ── Step 3 ────────────────────────────────────────────────────────────────────
\subsection{Step 3: [Title]}

% TODO: Describe what this step accomplishes.

\begin{proof}
    % TODO: Fill in the proof for Step 3.
    \[
        \text{(derivation goes here)}
    \]
\end{proof}

% ── Final Result ──────────────────────────────────────────────────────────────
\subsection{Final Result}

\begin{theorem}[Optimal Regression Decision Rule]
    % TODO: State the final theorem.
    Under assumptions \textbf{A1}--\textbf{A3}, the optimal decision rule is
    \[
        % TODO: Replace with the actual result.
        \hat{y}^{*} \;=\; \hat{f}(x).
    \]
\end{theorem}

\begin{proof}
    % TODO: Combine the steps above into a complete proof.
    Follows directly from Steps 1--3. \qed
\end{proof}

% ─────────────────────────────────────────────────────────────────────────────
\end{document}
